\chapter{Filter Correction, Injection, and Reset}
\label{chap:filter_correction_reset}

\section{Aiding Updates}

Low-rate aiding sensors produce timestamped measurements. The update path owns
the standard error-state EKF lifecycle \cite{groves2013,simon2006}:
\begin{itemize}
    \item selecting measurements eligible for the current update time;
    \item computing residuals and Jacobians against the nominal state;
    \item applying Kalman correction to the error state;
    \item injecting the correction into the nominal state;
    \item resetting the error-state mean and covariance coordinates.
\end{itemize}

\section{Error-State EKF Correction}

For a measurement model with predicted measurement $\hat{\mathbf{z}}$,
measurement $\tilde{\mathbf{z}}$, Jacobian $H$, and covariance $R$, use the
residual convention
\begin{equation}
\mathbf{r}
=
\tilde{\mathbf{z}}
-
\hat{\mathbf{z}}.
\label{eq:v1_measurement_residual}
\end{equation}
The innovation covariance and Kalman gain are
\begin{align}
S
&=
H P^{-} H^T + R,
\label{eq:v1_innovation_covariance}\\
K
&=
P^{-} H^T S^{-1}.
\label{eq:v1_kalman_gain}
\end{align}
The correction applied to the error state is
\begin{equation}
\delta\hat{\xb}^{+}
=
K\mathbf{r}.
\label{eq:v1_error_state_correction}
\end{equation}
The posterior covariance before nominal-state injection and reset is computed
with the Joseph form:
\begin{equation}
\boxed{
P^{+}
=
\left(\I-KH\right)P^{-}\left(\I-KH\right)^T
+
K R K^T.
}
\label{eq:v1_joseph_covariance_update}
\end{equation}

\section{Chi-Square Innovation Acceptance}
\label{sec:v1_chi_square_innovation_acceptance}

Innovation acceptance is a generic statistical filter capability, independent
of the sensor or observation model that produced the residual. For a
measurement model of dimension $m$, define the normalized innovation squared
(NIS)
\begin{equation}
\epsilon
=
\mathbf{r}^T S^{-1}\mathbf{r}.
\label{eq:v1_nis}
\end{equation}
Under the linearized Gaussian measurement hypothesis,
$\epsilon\sim\chi^2_m$. For a configured cumulative acceptance probability
$p_{\mathrm{acc}}$, with $0<p_{\mathrm{acc}}<1$, define the threshold
\begin{equation}
\gamma
=
F_{\chi^2_m}^{-1}\!\left(p_{\mathrm{acc}}\right).
\label{eq:v1_chi_square_gate_threshold}
\end{equation}
The degrees of freedom $m$ are derived from the selected observation model;
they are not an independent runtime input. The implementation likewise
computes $\gamma$ from the configured probability rather than accepting a raw
threshold literal.

Let $g_{\mathrm{enabled}}$ indicate whether statistical gating is enabled for
the observation. The complete acceptance condition is
\begin{equation}
\boxed{
\operatorname{accept}
=
\operatorname{valid}(S)
\land
\operatorname{finite}(\epsilon)
\land
\left(\epsilon\geq 0\right)
\land
\left(\neg g_{\mathrm{enabled}}\lor\epsilon\leq\gamma\right).
}
\label{eq:v1_chi_square_gate_acceptance}
\end{equation}
Disabling the statistical gate removes only the threshold comparison. A
numerically invalid innovation covariance or NIS is always rejected. The
decision is evaluated before persistent filter mutation: only an accepted
observation may accumulate an error-state correction or replace the covariance
used by a subsequent observation.

Every processed observation must retain diagnostics sufficient to reproduce
the decision: timestamp, acceptance status, innovation-covariance validity,
NIS, gate-enabled status, configured probability, model-derived degrees of
freedom, and computed threshold.

\section{Nominal-State Injection}

After an aiding update estimates $\delta\hat{\xb}^{+}$, apply translational and
bias corrections additively:
\begin{align}
\hat{\pb}_{eb}^{e,+} &\leftarrow \hat{\pb}_{eb}^{e,-} + \delta\hat{\pb}^{e},\\
\hat{\vb}_{eb}^{e,+} &\leftarrow \hat{\vb}_{eb}^{e,-} + \delta\hat{\vb}^{e},\\
\hat{\bb}_g^{+} &\leftarrow \hat{\bb}_g^{-} + \delta\hat{\bb}_g,\\
\hat{\bb}_a^{+} &\leftarrow \hat{\bb}_a^{-} + \delta\hat{\bb}_a.
\end{align}

For attitude, convert the small-angle correction into a scalar-first quaternion
correction:
\begin{equation}
\delta\hat{\quat}_{\theta}
\approx
\begin{bmatrix}
1\\
\frac{1}{2}\delta\hat{\boldsymbol{\theta}}
\end{bmatrix}.
\end{equation}

The NavKit v1 attitude-error convention in
\cref{eq:v1_attitude_error_dcm_relation} implies a left-multiplicative
correction for the whole nominal body-to-ECEF quaternion:
\begin{equation}
\boxed{
\hat{\quat}_{b}^{e,+}
\leftarrow
\delta\hat{\quat}_{\theta}
\otimes
\hat{\quat}_{b}^{e,-}.
}
\label{eq:v1_quat_injection}
\end{equation}
Then normalize:
\begin{equation}
\hat{\quat}^{+}
\leftarrow
\frac{\hat{\quat}^{+}}{\norm{\hat{\quat}^{+}}}.
\end{equation}

\begin{warningbox}
The correction side and sign in
\cref{eq:v1_quat_injection} fix the meaning of
$\delta\boldsymbol{\theta}$ and match
\cref{eq:v1_attitude_error_dcm_relation}. If implementation changes quaternion
storage direction or perturbation side, the prediction, measurement Jacobians,
injection, reset Jacobian, and tests must be updated together.
\end{warningbox}

\section{Reset}

After injection, reset the error-state mean to zero. Because attitude-error
coordinates change after multiplicative injection, reset the covariance with a
reset Jacobian:
\begin{equation}
P^{+}
\leftarrow
G_{\text{reset}}P^{+}G_{\text{reset}}^T.
\label{eq:v1_reset_covariance_update}
\end{equation}

For small attitude corrections, the attitude reset block is
\begin{equation}
G_{\text{reset},\theta}
\approx
\I+\frac{1}{2}\crossmat{\delta\hat{\boldsymbol{\theta}}}.
\label{eq:v1_attitude_reset_block}
\end{equation}
For the v1 state ordering, the full reset Jacobian is
\begin{equation}
\boxed{
G_{\text{reset}}
=
\begin{bmatrix}
\I & \zero & \zero & \zero & \zero\\
\zero & \I & \zero & \zero & \zero\\
\zero & \zero & G_{\text{reset},\theta} & \zero & \zero\\
\zero & \zero & \zero & \I & \zero\\
\zero & \zero & \zero & \zero & \I
\end{bmatrix}.
}
\label{eq:v1_full_reset_jacobian}
\end{equation}
Resolving the attitude error in ECEF keeps the error coordinates tied to the
navigation frame convention, but multiplicative attitude injection still
changes the local perturbation coordinates. The reset operation and reset
Jacobian therefore remain part of the filter/injection policy contract, not a
hidden side effect of a measurement model. The positive skew sign follows from
the left-multiplicative, ECEF-resolved error convention:
\begin{equation}
\exp\!\left(\crossmat{\delta\boldsymbol{\theta}^{+}}\right)
=
\exp\!\left(\crossmat{\delta\boldsymbol{\theta}^{-}}\right)
\exp\!\left(-\crossmat{\delta\hat{\boldsymbol{\theta}}}\right).
\end{equation}
Linearizing this composition about
$\delta\boldsymbol{\theta}^{-}=\delta\hat{\boldsymbol{\theta}}$ produces
$\I+\tfrac{1}{2}\crossmat{\delta\hat{\boldsymbol{\theta}}}$. The familiar
negative-skew reset block applies to the opposite, right-multiplicative local
error convention and must not be mixed with
\cref{eq:v1_attitude_error_dcm_relation,eq:v1_quat_injection}. This
approximation follows the same small-rotation convention as Sola's error-state
quaternion treatment \cite{sola2018}.
